% संपूर्ण मराठी ग्रंथातील प्रकरण 42: क्रमवर्ती कलन
% हा संपादनयोग्य स्रोतखंड आहे. संकलनासाठी संपूर्ण स्रोतसंग्रहातील
% अक्षररूपे, पूर्वभाग, चिन्हव्याख्या आणि आकृती-स्रोत आवश्यक आहेत.
% मूळ: Open Logic Project; मजकूर आणि रूपांतर CC BY 4.0.
% यंत्रानुवाद, दुरुस्ती आणि तपासणी: OpenAI Codex —
% GPT-5.6 Sol आणि GPT-6 Sol, दोन्ही Ultra effort.
% दुरुस्ती-पुनरावलोकन: GPT-6.1 Sol, Ultra effort.
\chapter{क्रमवर्ती कलन}\label{mvl:seq:chap}
\section{परिचय}\label{mvl:seq:int:sec}
अभिजात तर्कशास्त्रासाठीचे क्रमवर्ती कलन ही कार्यक्षम आणि सोपी
निष्पत्ती-पद्धत आहे. बहुमूल्य तर्कशास्त्राची व्याख्या ससीम
सत्यतामूल्ये असलेल्या मॅट्रिक्सने केली असेल, म्हणजे $V$ ससीम असेल,
तर त्यासाठी क्रमवर्ती कलन देता येते. हे कसे करावे याची कल्पना
अभिजात प्रसंगातील क्रमवर्तींचे अर्थ आणि निगमन नियमांचे रूप पाहिल्यावर
मिळते.
आता पुढील क्रमवर्ती लक्षात घ्या:
\begin{align*}
    \varphi_1, \dots, \varphi_m & \Sequent \psi_1, \dots, \psi_n
\intertext{हिचा अर्थ पुढील सूत्र म्हणून लावता येतो:}
    (\varphi_1 \land \cdots \land \varphi_m) & \lif (\psi_1 \lor \cdots \lor
\psi_n)
\end{align*}
दुसऱ्या शब्दांत, क्रमवर्ती $\Gamma \Sequent \Delta$ ची
सत्य-मूल्यांकन~$\pAssign{v}$ ने तेव्हा आणि केवळ तेव्हाच
\emph{पूर्ति होते},
जेव्हा काही $\varphi \in \Gamma$ साठी $\pValue{v}(\varphi) = \False$ असते
किंवा काही $\varphi \in \Delta$ साठी $\pValue{v}(\varphi) = \True$ असते.
या अर्थलक्षणानुसार, आरंभीच्या क्रमवर्ती $\varphi \Sequent \varphi$ ची
नेहमी पूर्ति होते; कारण एकतर $\pValue v(\varphi) = \True$ असते
किंवा $\pValue v(\varphi) = \False$ असते.
बाजूची सूत्रे $\Gamma$, $\Delta$ वगळल्यास, $\Log{LK}$ मधील
सशर्त संयोजकाचे निगमन नियम असे आहेत:
\begin{defish}
    \Axiom$ \fCenter \varphi$
    \Axiom$ \psi \fCenter $
    \RightLabel{\LeftR{\lif}}
    \BinaryInf$ \varphi \lif \psi \fCenter $
    \DisplayProof
    \hfill
    \Axiom$ \varphi \fCenter \psi$
    \RightLabel{\RightR{\lif}}
    \UnaryInf$ \fCenter \varphi \lif \psi $
    \DisplayProof
\end{defish}
क्रमवर्तीच्या वरील चिन्हार्थलक्षी अर्थलक्षणाचा उपयोग केल्यास,
$\LeftR{\lif}$ नियम असे सांगतो की $\pValue v(\varphi) = \True$
आणि $\pValue v(\psi) = \False$ असतील, तर $\pValue v(\varphi \lif \psi) = \False$
असते. त्याचप्रमाणे $\RightR{\lif}$ नियम असे सांगतो की
$\pValue v(\varphi) = \False$ किंवा $\pValue v(\psi) = \True$
असेल, तर $\pValue v(\varphi \lif \psi) = \True$ असते. प्रत्यक्षात,
येथील ही दोन्ही सशर्त विधाने द्विशर्त विधाने आहेत.
$\LeftR{\land}$ आणि $\RightR{\lor}$ नियमांच्या प्रमाणित रूपात
त्यांना अनुरूप सशर्त विधाने द्विशर्त नसतील. पण या नियमांची
अशी पर्यायी रूपे आहेत की त्यात ती द्विशर्त ठरतात:
\begin{defish}
    \Axiom$\varphi, \psi, \Gamma \fCenter \Delta$
    \RightLabel{\LeftR{\land}}
    \UnaryInf$\varphi \land \psi, \Gamma \fCenter \Delta$
    \DisplayProof
    \hfill
    \Axiom$ \Gamma \fCenter \Delta, \varphi, \psi$
    \RightLabel{\RightR{\lor}}
    \UnaryInf$ \Gamma \fCenter \Delta, \varphi \lor \psi$
    \DisplayProof
\end{defish}
ही मूलभूत कल्पना $n$-मूल्य तर्कशास्त्राला लावल्यास, दोन जागांऐवजी
$n$ जागा असलेले क्रमवर्ती कलन मिळते; प्रत्येक सत्यतामूल्यासाठी
एक जागा असते. $V = \{\False, \Undef, \True\}$ असलेल्या त्रिमूल्य
तर्कशास्त्रात क्रमवर्ती ही $\Gamma \mid \Pi \mid \Delta$ या रूपाची
अभिव्यक्ती असते. तिची सत्य-मूल्यांकन~$\pAssign v$ ने तेव्हा आणि
केवळ तेव्हाच पूर्ति होते, जेव्हा काही $\varphi \in \Gamma$ साठी
$\pValue{v}(\varphi) = \False$ असते, किंवा काही $\varphi \in \Delta$ साठी
$\pValue{v}(\varphi) = \True$ असते, किंवा काही $\varphi \in \Pi$ साठी
$\pValue{v}(\varphi) = \Undef$ असते. त्यामुळे आरंभीच्या क्रमवर्ती
$\varphi \mid \varphi \mid \varphi$ ची नेहमी पूर्ति होते.
\section{नियम आणि निष्पत्ती}\label{mvl:seq:rul:sec}
पुढे $\Gamma, \Delta, \Pi, \Lambda$ या वाक्यांच्या
सांत क्रमिका दर्शवतात असे समजू.
\begin{defn}[क्रमवर्ती]
\emph{$n$-बाजूंची क्रमवर्ती} म्हणजे पुढील रूपाची अभिव्यक्ती:
\[
\Gamma_1 \nSequent \dots \nSequent \Gamma_n
\]
येथे प्रत्येक $\Gamma_i$ ही भाषा~$\Lang L$ मधील
वाक्यांची सांत, शक्यतो रिकामी, क्रमिका असते.
\end{defn}

\begin{defn}[आरंभीची क्रमवर्ती]
\emph{$n$-बाजूंची आरंभीची क्रमवर्ती} म्हणजे भाषेतील कोणत्याही
वाक्य~$\varphi$ साठी $\varphi \nSequent \dots \nSequent \varphi$
या रूपाची $n$-बाजूंची क्रमवर्ती.

भाषेत $0$-स्थानी संयोजक~$\star$, म्हणजे विधानीय स्थिरांक,
असेल तर $\dots \nSequent \star \nSequent \dots$ ही क्रमवर्तीही
आरंभीची मानतो. यात $\tf{\star} \in V$ शी संबंधित
सत्यतामूल्याच्या जागी $\star$ असते; इतर जागा रिकाम्या असतात.
\end{defn}

$n$-मूल्य तर्कशास्त्र~$\Log{L}$ मधील प्रत्येक संयोजकासाठी,
त्या संयोजकाला $\Log{L}$ मध्ये मिळू शकणाऱ्या प्रत्येक
सत्यतामूल्यासाठी एक तार्किक नियम असतो. $\Log{L}$ च्या
$n$-बाजूंच्या क्रमवर्ती कलनातील निष्पत्त्या ही क्रमवर्तींची
सांत वृक्षरूप मांडणी असते: सर्वांत वरच्या क्रमवर्ती आरंभीच्या असतात;
आणि एखादी क्रमवर्ती एक किंवा अधिक क्रमवर्तींच्या खाली असेल,
तर ती या कलनातील निगमन नियमाने त्यांच्यापासून योग्यरीत्या
निष्पन्न झालेली असावी लागते. यात संयोजकांचे तार्किक नियम
आणि पुढील विभागात दिलेले संरचनात्मक नियम दोन्ही येतात.

\begin{defn}[प्रमेये]
वाक्य~$\varphi$ हे $n$-मूल्य तर्कशास्त्र~$\Log{L}$ चे
\emph{प्रमेय} तेव्हा असते, जेव्हा $\Log{L}$ मधील प्रत्येक निर्दिष्ट
सत्यतामूल्याशी संबंधित जागेत $\varphi$ असलेल्या $n$-बाजूंच्या
क्रमवर्तीची निष्पत्ती असते; या क्रमवर्तीतील
इतर सर्व जागा रिकाम्या असाव्यात. $\varphi$ प्रमेय असेल तर
$\Proves[\Log{L}] \varphi$ आणि नसेल तर $\Proves/[\Log{L}] \varphi$ लिहितो.
\end{defn}

\begin{defn}[निष्पन्नता]
$n$-मूल्य तर्कशास्त्र~$\Log{L}$ मध्ये वाक्य~$\varphi$
हे वाक्यांच्या संच~$\Gamma$ पासून \emph{निष्पन्न करता येण्याजोगे}
असते, $\Gamma \Proves[\Log{L}] \varphi$, तेव्हा आणि केवळ तेव्हाच,
जेव्हा एखादा ससीम उपसंच~$\Gamma_0 \subseteq \Gamma$
आणि $\Gamma_0$ मधील वाक्यांची एखादी क्रमिका~$\Gamma_0'$
अशी असते की पुढील क्रमवर्तीची निष्पत्ती असते:
\[ \Lambda_1 \nSequent \dots \nSequent \Lambda_n \]
येथे जागा~$i$ निर्दिष्ट सत्यतामूल्याशी संबंधित असेल तर
$\Lambda_i$ हे $\varphi$ असते; अन्यथा ते $\Gamma_0'$ असते.
$\varphi$ हे $\Gamma$ पासून निष्पन्न करता येण्याजोगे नसेल तर
$\Gamma \Proves/ \varphi$ असे लिहितो.
\end{defn}

उदाहरणार्थ, $3$-मूल्य \L ukasiewicz तर्कशास्त्रासाठी
$3$-बाजूंचे क्रमवर्ती कलन आहे. $3$-बाजूंच्या क्रमवर्ती
$\Gamma \nSequent \Pi \nSequent \Delta$ मध्ये $\Gamma$
हे~$\False$ शी, $\Delta$ हे~$\True$ शी आणि $\Pi$
हे~$\Undef$ शी संबंधित असतात. स्वयंसिद्धके
$\varphi \nSequent \varphi \nSequent \varphi$ ही आहेत. फक्त $\True$
निर्दिष्ट असल्यामुळे, $\Gamma$ ही सांत क्रमिका असल्यास
$\Gamma \Proves[\LogLuk[3]] \varphi$
तेव्हा आणि केवळ तेव्हाच, जेव्हा $\Gamma \nSequent \Gamma \nSequent \varphi$
या क्रमवर्तीची निष्पत्ती असते. ($\Undef$ देखील निर्दिष्ट
असते, तर $\Gamma \nSequent \varphi \nSequent \varphi$ या क्रमवर्तीची
निष्पत्ती आवश्यक झाले असते.)
अनंत आधारसूत्रसंचासाठी, आधीच्या व्याख्येनुसार त्यातील सांत
उपसंचाच्या क्रमिकेचा वापर करायचा; अनंत क्रमवर्तीचा नाही.












\section{संरचनात्मक नियम}\label{mvl:seq:str:sec}

$n$-बाजूंच्या क्रमवर्ती कलनातील संरचनात्मक नियम
अभिजात प्रसंगाप्रमाणेच कार्य करतात; मात्र ते प्रत्येक जागा~$i$
साठी स्वतंत्रपणे लागू होतात.

\begin{defish}
\begin{center}
\AxiomC{$ \Gamma_1 \nSequent \dots \nSequent \phantom{\varphi,}\Gamma_i \nSequent \dots \nSequent \Gamma_n $}
\RightLabel{$\iR{\Weakening}{i}$}
\UnaryInfC{$\Gamma_1 \nSequent \dots \nSequent \varphi, \Gamma_i \nSequent \dots \nSequent \Gamma_n$}
\DisplayProof
\\[2ex]
\AxiomC{$ \Gamma_1 \nSequent \dots \nSequent \varphi, \varphi, \Gamma_i \nSequent \dots \nSequent \Gamma_n $}
\RightLabel{$\iR{\Contraction}{i}$}
\UnaryInfC{$\Gamma_1 \nSequent \dots \nSequent \phantom{\varphi,}\varphi, \Gamma_i \nSequent
\dots \nSequent \Gamma_n$}
\DisplayProof
\\[2ex]
\AxiomC{$ \Gamma_1 \nSequent \dots \nSequent \Gamma_i, \varphi, \psi, \Gamma_i' \nSequent \dots \nSequent \Gamma_n $}
\RightLabel{$\iR{\Exchange}{i}$}
\UnaryInfC{$\Gamma_1 \nSequent \dots \nSequent \Gamma_i, \psi, \varphi, \Gamma_i' \nSequent \dots
\nSequent \Gamma_n$}
\DisplayProof
\end{center}
\end{defish}
दुर्बलीकरण, आकुंचन आणि अदलाबदल यांच्या सलग निगमनांची मालिका
अनेकदा दुहेरी निगमन रेषांनी दर्शवली जाते.
\Cut{} नियमाची अनेक रूपे असतात: क्रमवर्तीतील भिन्न जागांच्या
प्रत्येक जोडी~$i \neq j$ साठी एक रूप:
\begin{defish}
\[
\AxiomC{$ \Gamma_1 \nSequent \dots \nSequent \varphi, \Gamma_i \nSequent \dots \nSequent \Gamma_n $}
\AxiomC{$ \Delta_1 \nSequent \dots \nSequent \varphi, \Delta_j \nSequent \dots \nSequent \Delta_n $}
\RightLabel{$\iR{\Cut}{i,j}$}
\BinaryInfC{$\Gamma_1,\Delta_1 \nSequent \dots \nSequent \Gamma_n, \Delta_n$}
\DisplayProof
\]
\end{defish}












\section{निवडक तर्कशास्त्रांसाठी विधानीय नियम}\label{mvl:seq:prl:sec}

एखाद्या संयोजकासाठीचे $n$-बाजूंच्या क्रमवर्ती कलनातील निगमन नियम
फक्त त्या संयोजकाच्या वैशिष्ट्यपूर्ण सत्यता-फलनावर अवलंबून असतात.
म्हणून भिन्न तर्कशास्त्रांत एखाद्या संयोजकाची व्याख्या एकाच
सत्यता-फलनाने केली असेल, तर त्या संयोजकाचे $n$-बाजूंचे
क्रमवर्ती नियम त्या तर्कशास्त्रांत सारखेच असतात.

\subsection{$\lnot$ साठीचे नियम}

$\lnot$ साठीचे पुढील नियम \L ukasiewicz आणि क्लिनी यांच्या
तर्कशास्त्रांना तसेच त्यांच्या प्रकारांना लागू होतात.

\begin{defish}
  \begin{center}
\AxiomC{$ \Gamma \nSequent \Pi \nSequent \Delta, \varphi $}
\RightLabel{\iR{\lnot}{\False}}
\UnaryInfC{$ \lnot \varphi, \Gamma \nSequent \Pi \nSequent \Delta$}
\DisplayProof
\\[2ex]
\AxiomC{$ \Gamma \nSequent \varphi, \Pi \nSequent \Delta $}
\RightLabel{\iR{\lnot}{\Undef}}
\UnaryInfC{$\Gamma \nSequent \lnot \varphi, \Pi \nSequent \Delta$}
\DisplayProof
\\[2ex]
\AxiomC{$\varphi, \Gamma \nSequent \Pi \nSequent \Delta$}
\RightLabel{\iR{\lnot}{\True}}
\UnaryInfC{$\Gamma \nSequent \Pi \nSequent \Delta,  \lnot \varphi$}
\DisplayProof
  \end{center}
\end{defish}
$\lnot$ साठीचे पुढील नियम G\"odel तर्कशास्त्राला लागू होतात.
\begin{defish}
\AxiomC{$ \Gamma \nSequent \varphi, \Pi \nSequent \Delta, \varphi $}
\RightLabel{\iR{\lnot}{\False}[\LogGod]}
\UnaryInfC{$ \lnot \varphi, \Gamma \nSequent \Pi \nSequent \Delta$}
\DisplayProof
\hfill
\AxiomC{$\varphi, \Gamma \nSequent \Pi \nSequent \Delta$}
\RightLabel{\iR{\lnot}{\True}[\LogGod]}
\UnaryInfC{$\Gamma \nSequent \Pi \nSequent \Delta,  \lnot \varphi$}
\DisplayProof
\hfill
\end{defish}
(G\"odel तर्कशास्त्रात $\lnot \varphi$ ला~$\Undef$ हे मूल्य
कधीही मिळत नाही; म्हणून मधल्या जागेसाठी नियम नाही.)
\subsection{$\land$ साठीचे नियम}
\L ukasiewicz, प्रबळ क्लिनी आणि G\"odel यांच्या
तर्कशास्त्रांतील $\land$ साठीचे नियम पुढीलप्रमाणे आहेत.
\begin{defish}
\begin{center}
  \AxiomC{$\varphi, \psi, \Gamma \nSequent \Pi \nSequent \Delta$}
  \RightLabel{$\iR{\land}{\False}$}
  \UnaryInfC{$\varphi \land \psi, \Gamma \nSequent \Pi \nSequent \Delta$}
  \DisplayProof
\\[2ex]
  \AxiomC{$\Gamma \nSequent \varphi, \Pi \nSequent \varphi, \Delta$}
  \AxiomC{$\Gamma \nSequent \psi, \Pi \nSequent \psi, \Delta$}
  \AxiomC{$\Gamma \nSequent \varphi, \psi, \Pi \nSequent \Delta$}
  \RightLabel{$\iR\land\Undef$}
  \TrinaryInfC{$\Gamma \nSequent \varphi \land \psi, \Pi \nSequent \Delta$}
  \DisplayProof
  \\[2ex]
  \AxiomC{$\Gamma \nSequent \Pi \nSequent \Delta, \varphi$}
  \AxiomC{$\Gamma \nSequent \Pi \nSequent \Delta, \psi$}
  \RightLabel{$\iR\land\True$}
  \BinaryInfC{$\Gamma \nSequent \Pi \nSequent \Delta, \varphi \land \psi$}
  \DisplayProof
\end{center}
\end{defish}
\subsection{$\lor$ साठीचे नियम}
\L ukasiewicz, प्रबळ क्लिनी आणि G\"odel यांच्या
तर्कशास्त्रांतील $\lor$ साठीचे नियम पुढीलप्रमाणे आहेत.
\begin{defish}
\begin{center}
 \AxiomC{$\varphi, \Gamma \nSequent \Pi \nSequent \Delta$}
  \AxiomC{$\psi, \Gamma \nSequent \Pi \nSequent \Delta$}
  \RightLabel{$\iR\lor\False$}
  \BinaryInfC{$\varphi \lor \psi, \Gamma \nSequent \Pi \nSequent \Delta$}
  \DisplayProof
\\[2ex]
  \AxiomC{$\varphi, \Gamma \nSequent \varphi, \Pi \nSequent \Delta$}
  \AxiomC{$\psi, \Gamma \nSequent \psi, \Pi \nSequent \Delta$}
  \AxiomC{$\Gamma \nSequent \varphi, \psi, \Pi \nSequent \Delta$}
  \RightLabel{$\iR\lor\Undef$}
  \TrinaryInfC{$\Gamma \nSequent \varphi \lor \psi, \Pi \nSequent \Delta$}
  \DisplayProof
  \\[2ex]
  \AxiomC{$\Gamma \nSequent \Pi \nSequent \Delta, \varphi, \psi$}
  \RightLabel{$\iR\lor\True$}
  \UnaryInfC{$\Gamma \nSequent \Pi \nSequent \Delta, \varphi \lor \psi$}
  \DisplayProof
\end{center}
\end{defish}
\subsection{$\lif$ साठीचे नियम}
\L ukasiewicz तर्कशास्त्रातील $\lif$ साठीचे नियम पुढीलप्रमाणे आहेत.
\begin{defish}
\begin{center}
  \AxiomC{$\Gamma \nSequent \Pi \nSequent \Delta, \varphi$}
  \AxiomC{$\psi, \Gamma \nSequent \Pi \nSequent \Delta$}
  \RightLabel{$\iR\lif\False[\LogLuk[3]]$}
  \BinaryInfC{$\varphi \lif \psi, \Gamma \nSequent \Pi \nSequent \Delta$}
  \DisplayProof
  \\[2ex]
  \AxiomC{$\Gamma \nSequent \varphi, \psi, \Pi \nSequent \Delta$}
  \AxiomC{$\psi, \Gamma \nSequent \Pi \nSequent \Delta, \varphi$}
  \RightLabel{$\iR\lif\Undef[\LogLuk[3]]$}
  \BinaryInfC{$\Gamma \nSequent \varphi \lif \psi, \Pi \nSequent \Delta$}
  \DisplayProof
\\[2ex]
  \AxiomC{$\varphi, \Gamma \nSequent \psi, \Pi \nSequent \Delta, \psi$}
  \AxiomC{$\varphi, \Gamma \nSequent \varphi, \Pi \nSequent \Delta, \psi$}
  \RightLabel{$\iR{\lif}{\True}[\LogLuk[3]]$}
  \BinaryInfC{$\Gamma \nSequent \Pi \nSequent \Delta, \varphi \lif \psi$}
  \DisplayProof
\end{center}
\end{defish}
प्रबळ क्लिनी तर्कशास्त्रातील $\lif$ साठीचे नियम पुढीलप्रमाणे आहेत.
\begin{defish}
\begin{center}
  \AxiomC{$\Gamma \nSequent \Pi \nSequent \Delta, \varphi$}
  \AxiomC{$\psi, \Gamma \nSequent \Pi \nSequent \Delta$}
  \RightLabel{$\iR\lif\False[\LogKs]$}
  \BinaryInfC{$\varphi \lif \psi, \Gamma \nSequent \Pi \nSequent \Delta$}
  \DisplayProof
  \\[2ex]
  \AxiomC{$\psi, \Gamma \nSequent \psi, \Pi \nSequent \Delta$}
  \AxiomC{$\Gamma \nSequent \varphi, \psi, \Pi \nSequent \Delta$}
  \AxiomC{$\Gamma \nSequent \varphi, \Pi \nSequent \Delta, \varphi$}
  \RightLabel{$\iR\lif\Undef[\LogKs]$}
  \TrinaryInfC{$\Gamma \nSequent \varphi \lif \psi, \Pi \nSequent \Delta$}
  \DisplayProof
\\[2ex]
  \AxiomC{$\varphi, \Gamma \nSequent \Pi \nSequent \Delta, \psi$}
  \RightLabel{$\iR{\lif}{\True}[\LogKs]$}
  \UnaryInfC{$\Gamma \nSequent \Pi \nSequent \Delta, \varphi \lif \psi$}
  \DisplayProof
\end{center}
\end{defish}
G\"odel तर्कशास्त्रातील $\lif$ साठीचे नियम पुढीलप्रमाणे आहेत.
\begin{defish}
\begin{center}
  \AxiomC{$\Gamma \nSequent \varphi, \Pi \nSequent \Delta, \varphi$}
  \AxiomC{$\psi, \Gamma \nSequent \Pi \nSequent \Delta$}
  \RightLabel{$\iR\lif\False[\LogGod[3]]$}
  \BinaryInfC{$\varphi \lif \psi, \Gamma \nSequent \Pi \nSequent \Delta$}
  \DisplayProof
  \\[2ex]
  \AxiomC{$\Gamma \nSequent \psi, \Pi \nSequent \Delta$}
  \AxiomC{$\Gamma \nSequent \Pi \nSequent \Delta, \varphi$}
  \RightLabel{$\iR\lif\Undef[\LogGod[3]]$}
  \BinaryInfC{$\Gamma \nSequent \varphi \lif \psi, \Pi \nSequent \Delta$}
  \DisplayProof
\\[2ex]
  \AxiomC{$\varphi, \Gamma \nSequent \psi, \Pi \nSequent \Delta, \psi$}
  \AxiomC{$\varphi, \Gamma \nSequent \varphi, \Pi \nSequent \Delta, \psi$}
  \RightLabel{$\iR{\lif}{\True}[\LogGod[3]]$}
  \BinaryInfC{$\Gamma \nSequent \Pi \nSequent \Delta, \varphi \lif \psi$}
  \DisplayProof
\end{center}
\end{defish}
\begin{sidewaysfigure}
  \begin{center}
  \AxiomC{$A \nSequent A \nSequent A$}
  \RightLabel{\iR \Weakening \True}
  \UnaryInfC{$A \nSequent A \nSequent B, A$}
  \RightLabel{\iR \Weakening \Undef}
  \UnaryInfC{$A \nSequent B, A \nSequent B, A$}
  \RightLabel{\iR \Weakening \Undef}
  \UnaryInfC{$A \nSequent A, B, A \nSequent B, A$}
  \AxiomC{$A \nSequent A \nSequent A$}
  \RightLabel{\iR \Weakening \True}
  \UnaryInfC{$A \nSequent A \nSequent A, A$}
  \RightLabel{\iR \Weakening \True}
  \UnaryInfC{$A \nSequent A \nSequent B, A, A$}
  \RightLabel{\iR \Weakening \False}
  \UnaryInfC{$B, A \nSequent A \nSequent B, A, A$}
  \RightLabel{$\iR\lif\Undef$}
  \BinaryInfC{$A \nSequent A \lif B, A \nSequent B, A$}
  \AxiomC{$B \nSequent B \nSequent B$}
  \RightLabel{\iR \Weakening \Undef}
  \UnaryInfC{$B \nSequent A, B \nSequent B$}
  \RightLabel{\iR \Exchange \Undef}
  \UnaryInfC{$B \nSequent B, A \nSequent B$}
  \RightLabel{\iR \Weakening \Undef}
  \UnaryInfC{$B \nSequent A, B, A \nSequent B$}
  \RightLabel{\iR \Weakening \False}
  \UnaryInfC{$A, B \nSequent A, B, A \nSequent B$}
  \RightLabel{\iR \Exchange \False}
  \UnaryInfC{$B, A \nSequent A, B, A \nSequent B$}
  \AxiomC{$A \nSequent A \nSequent A$}
  \RightLabel{\iR \Weakening \True}
  \UnaryInfC{$A \nSequent A \nSequent B, A$}
  \RightLabel{\iR \Weakening \False}
  \UnaryInfC{$B, A \nSequent A \nSequent B, A$}
  \RightLabel{\iR \Weakening \False}
  \UnaryInfC{$B, B, A \nSequent A \nSequent B, A$}
  \RightLabel{$\iR\lif\Undef$}
  \BinaryInfC{$B, A \nSequent A \lif B, A \nSequent B$}
  \RightLabel{$\iR\lif\False$}
  \BinaryInfC{$A \lif B, A \nSequent A \lif B, A \nSequent B$}
  \DisplayProof
  \end{center}
  \caption{$\LogLuk[3]$ मधील निष्पत्तीचे उदाहरण}
\end{sidewaysfigure}
\part{सामान्य मोडल तर्कशास्त्रे}\label{nml:part}
\begin{editorial}
  या भागात सामान्य मोडल तर्कशास्त्रांच्या अधिसिद्धांताचा विचार केला आहे.
  सध्या यात मूलभूत मोडल तर्कशास्त्राच्या अभिजात अनुरूपता
  सिद्धांतावरील Aldo Antonelli यांची टिपणे आहेत.
\end{editorial}
